\documentclass[../main]{subfiles}
\begin{document}
\chapter{Ley de Faraday-Lenz}
\img{images/05/faradaylenz.png}{Experimento de Faraday.}
\info{F.E.M media generada por la variación de flujo en el tiempo}{
En un circuito cerrado muy delgado en el cual varía el flujo magnético se induce una fuerza electromotriz (F.E.M.) proporcional a la variación temporal del flujo. 
\begin{equation*}
   F.E.M. = E_m = -N \times \frac{\Delta \Phi}{\Delta t} = -N \times \frac{\Phi_{final} - \Phi_{inicial}}{t_{final}-t_{inicial}}
\end{equation*}

siendo:

$E_m:$ Fuerza electromotriz media generada [Volt]

$N$: Número de espiras.

$\Phi_{inicial}$: Flujo inicial [Webber]

$\Phi_{final}$: Flujo final [Webber]

$t_{inicial}$:Tiempo inicial [segundos]

$t_{final}$:Tiempo inicial [segundos]

}

\actividad{Resolver los siguientes problemas.}

\begin{enumerate}
    \item Una espira es atravesada por un flujo de $\Phi_{incial}=0,0018Wb$ y $\Delta t=0,01s$ segundos mas tarde el flujo que lo atraviesa es de $\Phi_{final}= 0,006Wb$. Calcular el valor medio de la f.e.m. $E_m$.
    \item Un flujo magnético que varía $\Delta \Phi = 0,03Wb$ atraviesa una bobina generando una f.e.m. de $E_m=30V$. Calcular el número de espiras $N$ si transcurre en un tiempo $\Delta t=1s$.
    \item Un solenoide de $N=2000$ espiras, longitud de $l=50cm$, $\varnothing=4cm$ y núcleo de madera el cual tiene arrollado sobre él una bobina de $N=1500$ espiras. Calcular la f.e.m. inducida en la bobina cuando se anula en $\Delta t=0,1s$. La corriente es de $I=10A$ 
    \item Una bobina de $N=400$ espiras se encuentra dentro de un campo magnético y es travesado por un flujo de $\Phi=1,8\times10^{-14}Mx$ y se saca el campo en $\Delta t=0,05s$. Calcular el valor medio de la f.e.m. inducida.
    \item Calcular la variación del flujo magnético en una bobina de $N=3000$ espiras si en $\Delta t = 1s$ se induce en ella una f.e.m de $E_m=60V$.
    \item Una bobina de $N=2000$ espiras se encuentra dentro de un campo magnético y su sección está atravesada por un flujo de $\Phi=0,005Wb$. Calcular el tiempo para que se genere una f.e.m. de $E_m=10V$.
    \item Una espira circular de radio $r=2cm$ se encuentra colocada en su sección perpendicularmente a las líneas de un campo magnético de inducción $B=12000Gs$. Si gira colocando su sección paralelamente a las líneas de fuerza en $\Delta t=0,1s$. Calcular la f.e.m. $E_m$.
\end{enumerate}
\end{document}